\subsection*{直线与平面的位置关系（4）}

知识点：三垂线定理

三垂线定理: 平面上的一条直线和这个平面的一条斜线垂直的充要条件是它和这条斜线在\blkda[3]{平面内的投影}垂直.

A 组：

\begin{enumerate}

\item 平面 $ a $ 外的直线 $ a $ 与平面 $ a $ 所成的角是 $ \theta $ ，则 $ \theta $ 的取值范围是\blkda{$[0,\dfrac{\pi}{2}]$}.

\item 在正四面体 $ {ABCD} $ 中，直线 $ {BC} $ 与 $ {AD} $ 所成角的大小为\blkda{$\dfrac{\pi}{2}$}.

3、定点 $ A $ 和 $ B $ 都在平面 $ \alpha $ 内,定点 $ P \notin  \alpha ,{PB} \bot  \alpha ,C $ 是平面 $ \alpha $ 内异于 $ A $ 和 $ B $ 的动点，且 $ {PC}\bot {AC} $ ，则 $ \bigtriangleup  {ABC} $ 为
\dotfill{(\kh{B})}

\xx{锐角三角形}{直角三角形} {钝角三角形} {无法确定}



\item 如图，如果 $ {MC}\bot $ 菱形 $ {ABCD} $ 所在平面，那么 $ {MA} $ 与 $ {BD} $ 的位置关系是\dotfill{(\kh{C})}
\xx{平行} {垂直相交} {垂直但不相交} {相交但不垂直}
\begin{figure}[h]
    \centering
    \includegraphics[width=0.25\linewidth]{bodyxb//src/afh'auhf'aouhfa.png}
\end{figure}


\item 在空间四边形 $ {ABCD} $ 中， $ E $ ， $ F $ ， $ G $ ， $ H $ 分别是边 $ {AB} $ ， $ {BC} $ ， $ {CD} $ ， $ {DA} $ 的中点， 若四边形对角线 $ {AC} = {BD} = 2 $ ，对角线 $ {AC} $ 与 $ {BD} $ 所成的角为 $ \frac{\pi }{3} $ ，则 $ {FH} = $ \blkda{$1or\sqrt{3}$}.

\end{enumerate}


B 组:
\begin{enumerate}

\item 直线 $ m $ 和平面 $ a $ 所成角为 $ \frac{\pi }{6} $ ，则直线 $ m $ 和平面 $ a $ 内任意直线所成角的取值范围是\blkda{$[\dfrac{\pi}{6},\dfrac{\pi}{2}]$}.

\item 在《九章算术》中，将底面为矩形且有一条侧棱与底面垂直的四棱锥称为 “阳马”. 如图，四棱锥 $ P - {ABCD} $ 为阳马，侧棱 $ {PA}\bot $ 底面 $ {ABCD} $ ， $ {PA} = {AB} = {AD} $ ， $ E $ 为棱 $ {PA} $ 的中点，则直线 $ {CE} $ 与平面 $ {PAD} $ 所成角的余弦值为\blkda{$\dfrac{\sqrt{5}}{3}$}.
\begin{figure}[h]
    \centering
    \includegraphics[width=0.25\linewidth]{bodyxb//src/a'ofuh'aushdfa.png}
\end{figure}



\item 如图，在坡面 $ a $ 与水平面 $ \beta $ 所成二面角为 $ {60}^{ \circ  } $ 的山坡上，有段直线型道路 $ {AB} $ 与坡脚 $ l $ 成 $ {30}^{ \circ  } $ 的角,这段路直通山顶 $ \mathrm{A} $ ，已知此山高 $ {135}\sqrt{3} $ 米，若小李从 $ B $ 沿着这条路上山，并且行进速度为每分钟 30 米，那么小李到达山顶 A 需要的时间是\blkda{18}分钟.
\begin{figure}[h]
    \centering
    \includegraphics[width=0.25\linewidth]{askfgagfafasafa.png}
\end{figure}


\item 已知点 $ P $ 是 $ \bigtriangleup  {ABC} $ 所在平面外一点，且 $ {PA} = {PB} = {PC} $ ，则点 $ P $ 在平面 $ {ABC} $ 上的射影一定是 $ \bigtriangleup  {ABC} $ 的\blkda{外}心.


\item 下列说法正确的是\dotfill{(\kh{B})}

\xx{过平面外一点作这个平面的垂直平面是唯一的}
{过平面的一条斜线作这个平面的垂直平面是唯一的}
{过直线外一点作这直线的平行平面是唯一的}
{过直线外一点作这直线的垂线是唯一的}



\item 如图， $ {AB} $ 为 $ \odot  O $ 的直径， $ {PA} $ 垂直于 $ \odot  O $ 所在的平面， $ M $ 为圆周上任意一点， $ {AN}\bot {PM} $ ， $ N $ 为垂足.

\begin{enumerate}
    \item 求证： $ {AN}\bot $ 平面 $ {PBM} $ ；
    \item 若 $ {AQ}\bot {PB} $ ，垂足为 $ Q $ ，求证： $ {NQ}\bot {PB} $ .
\end{enumerate}

\begin{figure}[h]
    \centering
    \includegraphics[width=0.25\linewidth]{asfashf'oauhfasd.png}
\end{figure}


\begin{jieda}
    \begin{enumerate}
        \item \( \because \mathrm{{AB}} \) 为 \( \odot  \mathrm{O} \) 的直径, \( \therefore \mathrm{{AM}}\bot \mathrm{{BM}} \) .
又 \( \mathrm{{PA}} \bot \) 平面 \( \mathrm{{ABM}} \) , \( \therefore \mathrm{{PA}} \bot  \mathrm{{BM}} \) ,
又 \( \because \mathrm{{PA}} \cap  \mathrm{{AM}} = \mathrm{A},\therefore \mathrm{{BM}}\bot \) 平面 \( \mathrm{{PAM}} \) .
又 \( \mathrm{{AN}} \subset \) 平面 \( \mathrm{{PAM}},\therefore \mathrm{{BM}}\bot \mathrm{{AN}} \) .
又 \( \mathrm{{AN}} \bot  \mathrm{{PM}} \) ,且 \( \mathrm{{BM}} \cap  \mathrm{{PM}} = \mathrm{M} \) ,
\( \therefore \mathrm{{AN}} \bot \) 平面 \( \mathrm{{PBM}} \) .
    \item 由(1)知 \( \mathrm{{AN}}\bot \) 平面 \( \mathrm{{PBM}},\mathrm{{PB}} \subset \) 平面 \( \mathrm{{PBM}},\therefore \mathrm{{AN}}\bot \mathrm{{PB}} \) .
又 \( \because \mathrm{{AQ}}\bot \mathrm{{PB}},\mathrm{{AN}} \cap  \mathrm{{AQ}} = \mathrm{A} \) ,
\( \therefore \mathrm{{PB}}\bot \) 平面 \( \mathrm{{ANQ}} \) .又 \( \mathrm{{NQ}} \subset \) 平面 \( \mathrm{{ANQ}} \) .
\( \therefore \mathrm{{PB}}\bot \mathrm{{NQ}} \) .
    \end{enumerate}
\end{jieda}



\item 已知长方体 $ {ABCD} - {A}_{1}{B}_{1}{C}_{1}{D}_{1} $ 中， $ {AD} = {AB} = 4 $ ， $ A{A}_{1} = 2,E,F $ 分别是 $ {AB},{A}_{1}{D}_{1} $ 的中点.

\begin{enumerate}
    \item 求证：直线 $ {EF}// $ 平面 $ B{B}_{1}{D}_{1}D $ ；
    \item 求直线 $ {EF} $ 与平面 $ {BC}{C}_{1}{B}_{1} $ 所成角的正弦值.
\end{enumerate}

\begin{figure}[h]
    \centering
    \includegraphics[width=0.25\linewidth]{bodyxb//src/aufhga'uhsf'afa.png}
\end{figure}

\begin{jieda}
    \begin{enumerate}
        \item 取$BD_1$的中点G，连接FG和BG，可以证明$FE//BG$,从而证明直线 $ {EF}// $ 平面 $ B{B}_{1}{D}_{1}D $ .

        \item 取$B_1C_1$的中点M，可证明GB在平面$ {BC}{C}_{1}{B}_{1} $的投影为$BM$，因此$\angle{GBM}$即为所求线面角，$|GM|=2$,$|MB|=2\sqrt{2}$，因此$sin\angle{GBM}=\dfrac{\sqrt{3}}{3}$。
    \end{enumerate}
    \begin{figure}[h]
        \centering
        \includegraphics[width=0.25\linewidth]{ao'uhfaoushdfa'df.png}
    \end{figure}
\end{jieda}


\item 如图，在四棱锥 $ P - {ABCD} $ 中，底面 $ {ABCD} $ 是矩形.已知 $ {AB} = 3,{AD} = 2,{PA} = 2,{PD} = 2\sqrt{2},\angle {PAB} = {60}^{ \circ  } $ . 
\begin{enumerate}
    \item 证明： $ {AD}\bot $ 平面 $ {PAB} $ ；
    \item 求异面直线 $ {PC} $ 与 $ {AD} $ 所成角的大小.
\end{enumerate}

\begin{figure}[h]
    \centering
    \includegraphics[width=0.25\linewidth]{'asofuhaouihf'ashf.png}
\end{figure}


\begin{jieda}
    \begin{enumerate}
        \item 由于四边形 \( {ABCD} \) 是矩形,所以 \( {AD}\bot {AB} \) ,
由于 \( P{A}^{2} + A{D}^{2} = P{D}^{2} \) ,所以 \( {AD}\bot {PA} \) ,
由于 \( {AB} \cap  {PA} = A,{AB},{PA} \subset \) 平面 \( {PAB} \) ,所以 \( {AD}\bot \) 平面 \( {PAB} \) .
        \item 由于 \( {AD}//{BC} \) ,所以 \( {BC}\bot \) 平面 \( {PAB} \) ,
由于 \( {PB} \subset \) 平面 \( {PAB} \) ,所以 \( {BC} \bot  {PB} \) ,
所以 \( \angle {PCB} \) 是异面直线 \( {PC} \) 与 \( {AD} \) 所成角,
\( {PB} = \sqrt{{2}^{2} + {3}^{2} - 2 \times  2 \times  3 \times  \cos {60}^{ \circ  }} = \sqrt{7}, \)
所以 \( \tan \angle {PCB} = \frac{\sqrt{7}}{2} \) ,由于 \( \angle {PCB} \) 是锐角,所以 \( \angle {PCB} = \arctan \frac{\sqrt{7}}{2} \) .
    \end{enumerate}
\end{jieda}

\item 如图， $ \bigtriangleup  {{ACD}} $ 和 $ \bigtriangleup  {{BCD}} $ 都是边长为 2 的等边三角形，

$ {{AB}} = \sqrt{{6}},{{EB}} \bot $ 平面 $ {BCD} $ .

\begin{enumerate}
    \item 证明： $ {EB}// $ 平面 $ {ACD} $ ；
    \item 求直线 $ {EB} $ 到平面 $ {ACD} $ 的距离。
\end{enumerate}

\begin{figure}
    \centering
    \includegraphics[width=0.25\linewidth]{asfudha'uhfawdfa.png}
\end{figure}

\begin{jieda}
    \begin{enumerate}
        \item 由于四边形 \( {ABCD} \) 是矩形,所以 \( {AD}\bot {AB} \) ,
由于 \( P{A}^{2} + A{D}^{2} = P{D}^{2} \) ,所以 \( {AD}\bot {PA} \) ,
由于 \( {AB} \cap  {PA} = A,{AB},{PA} \subset \) 平面 \( {PAB} \) ,所以 \( {AD}\bot \) 平面 \( {PAB} \) .
        \item  由第一问的证明，线到面的距离即点B到平面ACD的距离.取CD的中点O，因为BO垂直CD和AO，因此$BO\perp ACD$,所以$d=|BO|=\sqrt{3}$
    \end{enumerate}

    \begin{figure}[h]
        \centering
        \includegraphics[width=0.25\linewidth]{‘faoushf'aouhfasdf.png}
    \end{figure}
\end{jieda}





\item 如图，已知 $ A{A}_{1} \bot $ 平面 $ {ABC},B{B}_{1}//A{A}_{1},{AB} = {AC} = 3 $ ， $ {BC} = {2\sqrt{5}},\;{A{A}_{1}} = \sqrt{7},\;{B{B}_{1}} = {2\sqrt{7}} $ ，点 $ E $ 和 $ F $ 分别为 $ {BC} $ 和 $ {A}_{1}C $ 的中点.

\begin{enumerate}
    \item 求证： $ {AE}\bot $ 平面 $ {BC}{B}_{1} $ ；
    \item 求直线 $ {A}_{1}{B}_{1} $ 与平面 $ {BC}{B}_{1} $ 所成角的大小.
\end{enumerate}


\begin{figure}[h]
    \centering
    \includegraphics[width=0.25\linewidth]{asufha'suodhfafa.png}
\end{figure}


\begin{jieda}
    \begin{enumerate}
        \item 因为$AE\perp BB_1$,$AE\perp BC$,因此$AE\perp BCB_1$
        
        \item 如图,取 \( B{B}_{1} \) 的中点 \( M \) 和 \( {B}_{1}C \) 的中点 \( N \) ,连
接 \( {A}_{1}M,{A}_{1}N,{NE} \) .
因为 \( N \) 和 \( E \) 分别为 \( {B}_{1}C \) 和 \( {BC} \) 的中点,所以 \( {NE}// \)  \( {B}_{1}B,{NE} = \frac{1}{2}{B}_{1}B \) ,故 \( {NE}//{A}_{1}A \) 且 \( {NE} = {A}_{1}A \) ,
所以四边形 \( {NEA}{A}_{1} \) 为平行四边形,所以 \( {A}_{1}N//{AE} \) 且 \( {A}_{1}N = {AE} \) .
因为 \( {AB} = {AC},E \) 为 \( {BC} \) 的中点,所以 \( {AE}\bot {BC},{A}_{1}N\bot {BC} \) .
因为 \( A{A}_{1} \bot \) 平面 \( {ABC},B{B}_{1}//A{A}_{1} \) ,则易证得 \( B{B}_{1} \bot \) 平面 \( {ABC} \) .
因为 \( {AE} \subset \) 平面 \( {ABC} \) ,所以 \( B{B}_{1} \bot  {AE},{A}_{1}N \bot  B{B}_{1} \) ,
又 \( B{B}_{1} \cap  {BC} = B \) ,所以 \( {A}_{1}N \bot \) 平面 \( {BC}{B}_{1} \) ,
从而 \( \angle {A}_{1}{B}_{1}N \) 为直线 \( {A}_{1}{B}_{1} \) 与平面 \( {BC}{B}_{1} \) 的夹角.
在 \( \bigtriangleup {ABC} \) 中,可得 \( {AE} = 2 \) ,所以 \( {A}_{1}N = {AE} = 2 \) .
因为 \( {BM}//{AA}_{1},{BM} = {AA}_{1} \) ,所以 \( {A}_{1}M//{AB},{A}_{1}M = {AB} \) ,又由
\( {AB} \bot  B{B}_{1} \) ,得 \( {A}_{1}M \bot  B{B}_{1} \) .
在 \( \mathrm{{Rt}}\bigtriangleup {A}_{1}M{B}_{1} \) 中,可得 \( {A}_{1}{B}_{1} = \sqrt{{B}_{1}{M}^{2} + {A}_{1}{M}^{2}} = 4 \) .
在 \( \mathrm{{Rt}}\bigtriangleup {A}_{1}N{B}_{1} \) 中, \( \sin \angle {A}_{1}{B}_{1}N = \frac{{A}_{1}N}{{A}_{1}{B}_{1}} = \frac{1}{2} \) ,因此 \( \angle {A}_{1}{B}_{1}N = {30}^{ \circ  } \) .
所以直线 \( {A}_{1}{B}_{1} \) 与平面 \( {BC}{B}_{1} \) 的夹角为 \( {30}^{ \circ  } \) .
    \end{enumerate}
    \begin{figure}[h]
        \centering
        \includegraphics[width=0.25\linewidth]{aoufuh'aouhf'aodfa.png}
    \end{figure}
\end{jieda}

\item 如图，已知 $ {P} $ 是平面 $ {A}{B}{C} $ 外一点， $ {P}{A} \bot $ 平面 $ {ABC},{AC} \bot  {BC} $ .

\begin{enumerate}
    \item 证明： $ {{BC}} \bot $ 平面 $ {{PAC}} $ ；
    \item 过 $ {A} $ 点作 $ {{AD}} $ 垂直 $ {{PC}} $ 于 $ {D} $ ，证明： $ {{AD}} \bot  {{PB}} $ ； 
    \item 若 $ {{BC}} = {1},{{AC}} = \sqrt{{3}} $ ，求点 $ {C} $ 到平面 $ {{PAB}} $的距离.
\end{enumerate}

\begin{figure}[h]
    \centering
    \includegraphics[width=0.25\linewidth]{afouh'aoushfadfa.png}
\end{figure}

\begin{jieda}
    \begin{enumerate}
        \item 因为 \( {PA}\bot \) 平面 \( {ABC} \) , \( {BC} \subset \) 平面 \( {ABC} \) ,所以 \( {PA}\bot {BC} \) ,
又 \( {AC}\bot {BC},{PA} \cap  {AC} = A,{PA},{AC} \subset \) 平面 \( {PAC} \) ,
所以 \( {BC} \bot \) 平面 \( {PAC} \) ;
        \item 因为 \( {BC}\bot \) 平面 \( {PAC},{AD} \subset \) 平面 \( {PAC} \) ,
所以 \( {BC}\bot {AD} \) ,又 \( {AD}\bot {PC},{PC} \cap  {BC} = C,{PC},{BC} \subset \) 平面 \( {PBC} \) ,
所以 \( {AD}\bot \) 平面 \( {PBC} \) ,又 \( {PB} \subset \) 平面 \( {PBC} \) ,所以 \( {AD}\bot {PB} \) ；
    \item 在平面 \( {ABC} \) 内过点 \( C \) 作 \( {CE}\bot {AB} \) 交 \( {AB} \) 于点 \( E \) ,
因为 \( {BC} = 1,{AC} = \sqrt{3},{AC}\bot {BC} \) ,所以 \( {AB} = \sqrt{{1}^{2} + {\left( \sqrt{3}\right) }^{2}} = 2 \) ’
所以 \( {CE} = \frac{{AC} \cdot  {BC}}{AB} = \frac{\sqrt{3}}{2} \) ,
又 \( {PA}\bot \) 平面 \( {ABC},{CE} \subset \) 平面 \( {ABC} \) ,所以 \( {PA}\bot {CE} \) ,
又 \( {PA} \cap  {AB} = A,{PA},{AB} \subset \) 平面 \( {PAB} \) ,所以 \( {CE}\bot \) 平面 \( {PAB} \) ,
所以点 \( C \) 到平面 \( {PAB} \) 的距离为 \( {CE} = \frac{\sqrt{3}}{2} \) .

    \begin{figure}[h]
        \centering
        \includegraphics[width=0.25\linewidth]{asdfhauhf'ads.png}
    \end{figure}
    \end{enumerate}
\end{jieda}



\end{enumerate}





\section*{平面与平面的位置关系}


\subsection*{平面与平面的位置关系（1）}

知识点：面面平行

A 组：
\begin{enumerate}

\item 判断下列命题是否正确：

    \begin{enumerate}
        \item 若一个平面内有两条直线平行于另一个平面，则这两个平面平行；\blkda{\ding{55}}

        \item 若一个平面内有无数条直线平行于另一个平面，则这两个平面平行； \blkda{\ding{55}}

        \item 若一个平面内任意一条直线平行于另一个平面，则这两个平面平行； \blkda{\ding{51}}

        \item 若一个平面内有两条相交直线，分别平行于另一个平面内的两条直线，则这两个平面平行; \blkda{\ding{51}}

        \item 若一个平面内一个角（锐角或钝角）的两边和另一个平面内的一个角的两边分别平行， 则这两个平面平行. \blkda{\ding{51}}

        
    \end{enumerate}

    \item 已知 $ a \subset  a $ ，则 $ a//\beta $ 是 $ a//\beta $ 的 \blkda[3]{充分不必要}条件

    \item 已知 $ m,n $ 为两条不同直线， $ a,\beta ,\gamma $ 为三个不同平面，则下列说法正确的是\dotfill{(\kh{D})}
    \xx{若 $ m//a,n \subset  a $ ，则 $ m//n $ } {若 $ m//a,a//\beta $ ，则 $ m//\beta $}
    {若 $ m//a,m//\beta $ ，则 $ a//\beta $} {若 $ a//\beta ,\beta //\gamma $ ，则 $ a//\gamma $}

    \item 过平面 $ a $ 外一条直线 $ a $ 与 $ a $ 平行的平面的个数是 \dotfill{(\kh{C})}

    \xx{只有一个} {至少有一个} {至多有一个} { 一个也没有}

    \item 经过平面外一点可以作\blkda{1}个平面平行于这个平面；可以作\blkda{无数}条直线平行于这个平面.
    
    \item 分别在两个平面内的三角形都有两条边与同一条直线垂直，则这两个三角所在面的位置关系是\blkda{平行}.

\end{enumerate}


B组：
\begin{enumerate}

    \item 下列条件中，可以判断平面 $ \alpha $ 与 $ \beta $ 平行的是 \dotfill{(\kh{D})}

    \xx{直线 $ a $ 与 $ a,\beta $ 相交且所成的角相等}
    {$ \alpha $ 内存在不共线的三点到 $ \beta $ 的距离相等}
    {l, $ {m} $ 是 $ \alpha $ 内两条直线,且 $ l//\beta ,{m}//\beta $}
    {$ l,m $ 是两条异面直线，且 $ l//\alpha ,m//\alpha ,l//\beta ,m//\beta $}

    \item 已知平面 $ \alpha // $ 平面 $ \beta $ ，直线 $ a\text{ 、 }b $ 分别与 $ \alpha \text{ 、 }\beta $ 所成的角相等，则直线 $ a\text{ 、 }b $ 的位置关系是\blkda[3]{相交、平行或异面}.

    \item 如图，平面 $ a// $ 平面 $ \beta $ ，自点 $ {O} $ 引三条直线分别交 $ \alpha $ 、 $ \beta $ 于点 $ {A} $ 、B、C 和点 $ {A}_{1} $ 、 $ {B}_{1} $ 、 $ {C}_{1} $ ，则 $ \bigtriangleup  {ABC} $ 与 $ \bigtriangleup  {A}_{1}{B}_{1}{C}_{1} $ 的关系是\blkda{相似}.

    \begin{center}
    
        \includegraphics[width=0.25\linewidth]{bodyxb//src/a'ouhsf'ouahs'fouhasodf.png}
    \end{center}

    
    \item 已知 $ {a}\text{ 、 }{\;b}\text{ 、 }{c} $ 为三条不重合的直线， $ a,\beta ,\gamma $ 为三个不重合的平面， 下列六个命题中正确的命题是\blkda{(1)(4)}.

    \begin{enumerate}
        \item $ a//c,b//c \Rightarrow  a//b;\; $ 
        \item $ a//\gamma ,b//\gamma  \Rightarrow  a//b $ ;
        \item  $ {c}//a,{c}//\beta  \Rightarrow  a//\beta ;\; $
        \item $ a//\gamma ,\beta //\gamma  \Rightarrow  a//\beta $ ;
        \item $ {a}//{c},{c}//a \Rightarrow  {a}//a; $ 
        \item $ {a}//\gamma ,a//\gamma  \Rightarrow  {a}//a $ .
    \end{enumerate}


    \item 已知平面 $ a//\beta $ ，点 $ A,C \in  a,B,D \in  \beta $ ，直线 $ {{AB}} $ 与 $ {{CD}} $ 交于点 $ {S} $ ，且 $ {{AS}} = 8,{{BS}} = 9 $ ， CD=34，则 SC 的长度是\blkda[3]{272或16}.

    \item 已知平面 $ a//\beta $ ，且 $ a,\beta $ 的距离为 $ {d} $ ， $ a \subset  a $ ，则在平面 $ \beta $ 中 \dotfill{(\kh{C})}
    \xx{有且只有一条直线与 $ {a} $ 的距离为 $ {d} $}
    {所有直线与 $ {a} $ 的距离都等于 $ {d} $}
    {有无数条直线与 $ a $ 的距离等于 $ d $}
    {所有直线与 $ a $ 距离都不等于 $ d $}

    \item 已知 AB、CD 是夹在两平行平面 $ a,\beta $ 之间的两条线段， $ {AB}\bot {CD},{AB} = 2,{AB} $ 与平面 $ a $ 成 $ {30}^{ \circ  } $ 角，则线段 CD 的取值范围是 \dotfill{(\kh{B})}

    \xx{$ \left( {\dfrac{2\sqrt{3}}{3},2\sqrt{3}}\right) $} {$ \left\lbrack  {\dfrac{2\sqrt{3}}{3}, + \infty }\right) $} {$ \left( {1,\dfrac{2\sqrt{3}}{3}}\right) $ }{$ \lbrack 1, + \infty ) $}
    
    \item 已知 $ {P} $ 是 $ \bigtriangleup  {ABC} $ 所在平面外一点， $ {A}_{1} $ 、 $ {B}_{1} $ 、 $ {C}_{1} $ 分别是 $ \bigtriangleup  {PBC} $ 、 $ \bigtriangleup  {PCA} $ 、 $ \bigtriangleup  {PAB} $ 的重心. 
    \begin{enumerate}
        \item 求证: 平面 $ {A}_{1}{B}_{1}{C}_{1}// $ 平面 $ {{ABC}} $ ;
        \item 求 $ {S}_{\bigtriangleup {A}_{1}{B}_{1}{C}_{1}} : {S}_{\bigtriangleup {ABC}}$.
    \end{enumerate}

    \begin{jieda}
    \begin{enumerate}

        \item 连接 $ P{A}^{\prime } $ 、 $ P{C}^{\prime \prime } $ ,并延长分别交BC、AB于点M,N,连接MN. $ \because {A}^{\prime }\text{、}{C}^{\prime \prime } $ 分别是 $ \bigtriangleup  {{PBC}} $ 、 $ \bigtriangleup  {{PAB}} $ 的重心,
$ \therefore P{A}^{\prime } = \dfrac{2}{3}{PM},P{C}^{\prime \prime } = \dfrac{2}{3}{PN} $ ,
$ \therefore {A}^{\prime }{C}^{\prime \prime }\parallel {MN} $ .
$ \because {A}^{\prime }{C}^{\prime \prime } ⊄ $ 平面 $ {{ABC}},{MN} \subset $ 平面 $ {{ABC}} $ ,
$ \therefore {A}^{\prime }{C}^{\prime \prime }\parallel $ 平面 $ {{ABC}} $ .
同理, $ {A}^{\prime }{B}^{\prime \prime }\parallel $ 平面 $ {{ABC}} $ .
$ \because {A}^{\prime }{C}^{\prime \prime } \cap  {A}^{\prime }{B}^{\prime } = {A}^{\prime } $ ,且 $ {A}^{\prime }{C}^{\prime \prime }\text{、}{A}^{\prime }{B}^{\prime } \subset $ 平面 $ {A}^{\prime }{B}^{\prime }{C}^{\prime \prime } $ ,
$ \therefore $ 平面 $ {A}^{\prime }{B}^{\prime }{C}^{\prime \prime }\parallel $ 平面 $ {ABC} $ .
        \begin{center}
            \includegraphics[width=0.25\linewidth]{bodyxb//src/auhf'aoshfasd.png}
        \end{center}

\item 由(1)知 $ {A}^{\prime }{C}^{\prime \prime } = \dfrac{2}{3}{MN},{A}^{\prime }{C}^{\prime \prime }\parallel {MN} $ .
$ \because {MN} = \dfrac{1}{2}{AC},{MN}\parallel {AC} $ ,
$ \therefore {A}^{\prime }{C}^{\prime \prime } = \dfrac{1}{3}{AC},{A}^{\prime }{C}^{\prime \prime }\parallel {AC} $ .
同理可得： $ {A}^{{}^{\prime }}{B}^{{}^{\prime }} = \dfrac{1}{3}{AB} $ , $ {A}^{{}^{\prime }}{B}^{{}^{\prime }} \parallel  {AB} $ , $ {B}^{{}^{\prime }}{C}^{{}^{\prime \prime }} = \dfrac{1}{3}{BC} $ , $ {B}^{{}^{\prime }}{C}^{{}^{\prime \prime }} \parallel  {BC} $ ,
则 $ \Delta {A}^{\prime }{B}^{\prime }{C}^{\prime \prime } \sim  {\Delta ABC} $ .
故 $ \Delta {A}^{\prime }{B}^{\prime }{C}^{\prime \prime } $ 与 $ \bigtriangleup  {{ABC}} $ 的面积之比为 $ \dfrac{{S}_{\Delta {A}^{\prime }{B}^{\prime }{C}^{\prime \prime }}}{{S}_{\Delta ABC}} = \dfrac{1}{9} $ .
    \end{enumerate}
    \end{jieda}
    \vspace{3cm}



    \item 如图，在四棱锥 $ C - {ABED} $ 中，四边形 $ {ABED} $ 是正方形，点 $ G $ ， $ F $ 分别是线段 $ {EC} $ ， $ {BD} $ 的中点。

    \begin{enumerate}
        \item 求证： $ {GF}// $ 平面 $ {ABC} $ ；
        \item 若 $ H $ 是线段 $ {BC} $ 的中点，证明：平面 $ {GFH}// $ 平面 $ {ACD} $ .
    \end{enumerate}
    
\begin{center}

    \includegraphics[width=0.25\linewidth]{auufh'aoushf'aouhfa.png}
\end{center}

\begin{jieda}
    \begin{enumerate}
        \item 证明：由四边形 $ {ABED} $ 为正方形可知,
连接 $ {AE} $ 必与 $ {BD} $ 相交于中点 $ F $ ,故 $ {GF} \parallel  {AC} $ ,
$ \because {GF} ⊄ $ 面 $ {ABC},{AC} \subset $ 平面 $ {ABC} $ ,
$ \therefore {GF} \parallel $ 面 $ {ABC} $ .

        \begin{center}
        
            \includegraphics[width=0.25\linewidth]{bodyxb//src/au;ghf;aighf'aisfa.png}
        \end{center}

        \item 由点 $ G,H $ 分别为 $ {CE},{CB} $ 中点可得： $ {GH} \parallel  {EB} \parallel  {AD} $ ,
$ \because {GH} ⊄ $ 面 $ {ACD},{AD} \subset $ 平面 $ {ACD},\therefore {GH} \parallel $ 平面 $ {ACD} $ ,
又由( 1 )可知, $ {GF} \parallel $ 平面 $ {ACD} $ ,
且 $ {GH} \cap  {GF} = G,{GH},{GF} \subset $ 平面 $ {GFH} $ ,
故平面 $ {GFH}\parallel $ 平面 $ {ACD} $ .
    \end{enumerate}
\end{jieda}


\item 如图，在四棱锥 $ P - {ABCD} $ 中， $ {PA} = 3 $ ， $ {AB} = 2 $ ，四边形 $ {ABCD} $ 为菱形， $ \angle {ABC} = \dfrac{\pi }{3} $ ， $ {PA}\bot $ 平面 $ {ABCD} $ ,E,F,Q 分别是 BC,PC,PD 的中点. 证明: 平面 $ {EFQ}// $ 平面 $ {PAB} $ ;

\begin{jieda}
    因为四边形 $ {ABCD} $ 为菱形,所以 $ {AB}//{CD} $ ,
又 $ E,F,Q $ 分别是 $ {BC},{PC},{PD} $ 的中点,
所以 $ {FQ}//{CD},{EF}//{PB} $ ,故 $ {FQ}//{AB} $ ,
因为 $ {EF} ⊄ $ 平面 $ {PAB},{PB} \subset $ 平面 $ {PAB} $ ,
所以 $ {EF}// $ 平面 $ {PAB} $ ,同理可得 $ {FQ}// $ 平面 $ {PAB} $ .
因为 $ {EF} \cap  {FQ} = F,{EF},{FQ} \subset $ 平面 $ {EFQ} $ ,
所以平面 $ {EFQ}// $ 平面 $ {PAB} $ .
\end{jieda}


\item 如图，已知 $ P $ 是平行四边形 $ {ABCD} $ 所在平面外一点， $ M $ 、 $ N $ 分别是 $ {AB} $ 、 $ {PC} $ 的三等分点 $ \left( {M\text{ 靠近 }B,N\text{ 靠近 }C}\right) $ ；

\begin{enumerate}
    \item 求证： $ {MN}// $ 平面 $ {PAD} $ .
    \item 在 $ {PB} $ 上确定一点 $ Q $ ，使平面 $ {MNQ}// $ 平面 $ {PAD} $ ，并证明.
\end{enumerate}

\begin{center}

    \includegraphics[width=0.25\linewidth]{bodyxb//src/ahkfa'khfg'afsd.png}
\end{center}

\begin{jieda}
    \begin{enumerate}
        \item 证明：过点 $ N $ 作 $ {NE}//{CD} $ ,交 $ {PD} $ 于点 $ E $ ,连接 $ {AE} $ ,
因为 $ N $ 为 $ {PC} $ 的三等分点,可得 $ {NE} = \dfrac{2}{3}{CD} $ ,
又因为 $ M $ 为 $ {AB} $ 的三等分点,可得 $ {AM} = \dfrac{2}{3}{AB} $ ,
因为 $ {AB}//{CD} $ 且 $ {AB} = {CD} $ ,所以 $ {AM}//{NE} $ 且 $ {AM} = {NE} $ ,
所以四边形 $ {AMNE} $ 为平行四边形,所以 $ {MN}//{AE} $ ,
又由 $ {MN} ⊄ $ 平面 $ {PAD},{AE} \subset $ 平面 $ {PAD} $ ,所以 $ {MN}// $ 平面 $ {PAD} $ .
        \item 证明：取 $ {PB} $ 取一点 $ Q $ ,使得 $ {BQ} = \dfrac{1}{3}{BP} $ ,即点 $ Q $ 为 $ {PB} $ 上靠近点 $ B $ 的三等点,
在 $ {\Delta PAB} $ 中,因为 $ M,Q $ 分别为 $ {AB},{PB} $ 的三等分点,可得 $ \dfrac{MB}{AB} = \dfrac{BQ}{BP} $ ,所以 $ {MQ}//{PA} $ ,
因为 $ {MQ} ⊄ $ 平面 $ {PAD},{PA} \subset $ 平面 $ {PAD} $ ,所以 $ {MQ}// $ 平面 $ {PAD} $ ;
又由 (1) 知 $ {MN}// $ 平面 $ {PAD} $ ,且 $ {MN} \cap  {MQ} = M,{MN},{MQ} \subset $ 平面 $ {MNQ} $ ,
所以平面 $ {MNQ}// $ 平面 $ {PAD} $ ,
即当点 $ Q $ 为 $ {PB} $ 上靠近点 $ B $ 的三等点时,能使得平面 $ {MNQ}// $ 平面 $ {PAD} $ .

\begin{center}

    \includegraphics[width=0.25\linewidth]{bodyxb//src/aufhaosufha'ouhdf.png}
\end{center}
    \end{enumerate}
\end{jieda}



\item 如图，在四棱锥 $ P - {ABCD} $ 中，四边形 $ {ABCD} $ 是正方形， $ {PA} = {PC} $ ，E 为侧棱 PD 上的点，且 $ {PE} = {3ED} $ .
\begin{enumerate}
    \item 证明： $ {PD} \bot  {AC} $ ；
    \item 在侧棱PC上是否存在一点 $ {F} $ ，使得 $ {BF}// $ 平面 $ {ACE} $ ？若存在,求 $ \dfrac{PC}{PF} $ 的值; 若不存在,请说明理由.

\end{enumerate}

\begin{center}

    \includegraphics[width=0.25\linewidth]{bodyxb//src/a;kjsfh'akhf'ashdfds.png}
\end{center}

\begin{jieda}
    \begin{enumerate}
        \item 设 $ {BD} $ 交 $ {AC} $ 于点 $ O $ ,连接 $ {OP} $ ,正方形 $ {ABCD} $ 中,则 $ {AO} = {OC} $ , $ {BD}\bot {AC} $ ,
又 $ {PA} = {PC} $ ,则 $ {PO}\bot {AC} $ ,又 $ {PO} \cap  {BD} = O,{PO},{BD} \subset $ 平面 $ {PBD} $ ,
因此 $ {AC} \bot $ 平面 $ {PBD} $ ,而 $ {PD} \subset $ 平面 $ {PBD} $ ,
所以 $ {PD} \bot  {AC} $ .
        
        \begin{center}
        
            \includegraphics[width=0.25\linewidth]{bodyxb//src/;yafga;iygf;afawsfa.png}
        \end{center}

        \item 侧棱 $ {PC} $ 上存在一点 $ F $ ,满足条件,
证明如下：如图,正方形 $ {ABCD} $ 中, $ {BO} = {OD} $ ,
在线段 $ {PE} $ 取一点 $ G $ ,使得 $ {GE} = {ED} $ ,由 $ {PE} = {3ED} $ ,得 $ \dfrac{PG}{PE} = \dfrac{2}{3} $ ,
连接 $ {BG},{FG} $ ,则 $ {BG}//{OE} $ ,而 $ {OE} \subset $ 平面 $ {ACE},{BG} ⊄ $ 平面 $ {ACE} $ ,
$ {则BG}// $ 平面 $ {ACE} $ ,由 $ {BF}// $ 平面 $ {ACE},{BG} \cap  {BF} = B,{BG},{BF} \subset $ 平面 $ {BGF} $ ,
得平面 $ {BGF}// $ 平面 $ {ACE} $ ,而平面 $ {PCD} \cap $ 平面 $ {ACE} = {EC} $ ,平面 $ {PCD} \cap $ 平面 $ {BGF} = {GF} $ ,
于是 $ {GF}//{EC},\dfrac{PF}{PC} = \dfrac{PG}{PE} = \dfrac{2}{3} $ ,
所以 $ \dfrac{PC}{PF} = \dfrac{3}{2} $ .


        \begin{center}
        
            \includegraphics[width=0.25\linewidth]{bodyxb//src/asufha'ufh'auwsdfa.png}
        \end{center}

    \end{enumerate}
\end{jieda}


\item 已知 $ \alpha//\beta //\gamma $ ，A、D $ \in  a $ ，C、F $ \in  \gamma $ ，AC、DF 分别交平面 $ \beta $ 于 $ {B} $ 、E. 
\begin{enumerate}
    \item 求证: $ \dfrac{AB}{BC} = \dfrac{DE}{EF} $ ;
    \item 设 $ {{AF}} $ 交 $ \beta $ 于 $ {M} $ ,若 $ \dfrac{AB}{BC} = 1 $ ，异面直线 $ {{AD}} $ 与 $ {{CF}} $ 所成角为 $ {60}^{ \circ  },{{AD}} = 2,{{CF}} = 4 $ ，求 $ {S}_{\bigtriangleup {BEM}} $ .
\end{enumerate}

\begin{center}

    \includegraphics[width=0.25\linewidth]{bodyxb//src/fa'ufha'ouhfa'odfs.png}
\end{center}

\begin{enumerate}
    \item 连接AF交平面β于G,连接AD,BE,CF,BG,EG.
由 $ \beta //\gamma $ ,平面 $ {{ACF}} \cap  \beta  = {{BG}} $ ,平面 $ {{ACF}} \cap  \gamma  = {{CF}} $ ,所以 $ {{BG}}//{{CF}} $ ,则 $ \dfrac{{{AB}}}{{{BC}}} = \dfrac{{{AG}}}{{{GF}}} $ ,
同理,由 $ \alpha //\beta $ ,可得 $ {{GE}}//{{AD}} $ ,则 $ \dfrac{{{DE}}}{{{EF}}} = \dfrac{{{AG}}}{{{GF}}} $ .
所以 $ \dfrac{{{AB}}}{{{BC}}} = \dfrac{{{DE}}}{{{EF}}} $ ；
    \item 因为BG//CF,GE//AD,所以 $ \angle $ BGE(或其补角)就是直线AD与CF所成的角.
因为 $ \dfrac{{{AB}}}{{{AC}}} = \dfrac{{{BG}}}{{{CF}}} = \dfrac{1}{2},\dfrac{{{GF}}}{{{AF}}} = \dfrac{{{GE}}}{{{AD}}} = \dfrac{1}{2} $ ,所以 $ {{BG}} = 2,{{GE}} = 1 $ ,
又 $ {{BE}} = \sqrt{7},{{CF}} = 4 $ ,
由余弦定理可得 $ \cos \angle {{BGE}} = \dfrac{1 + 4 - 7}{4} =  - \dfrac{1}{2} $ ,得 $ \angle {{BGE}} = {120}^{ \circ  } $ .
即直线 $ {{AD}} $ 与 $ {{CF}} $ 所成的角为 $ {60}^{ \circ  } $ .

\begin{center}

    \includegraphics[width=0.25\linewidth]{bodyxb//src/aksufha'siuhfasd.png}
\end{center}

\end{enumerate}



\item 如图所示，在四棱锥 $ P - {ABCD} $ 中，底面 $ {ABCD} $ 为梯形， $ {BC}//{AD} $ ， $ {BC} = \dfrac{1}{2}{AD} $ ，面 $ {PBC} \cap $ 面 $ {PAD} = l $ ， $ E $ 是 $ {PD} $ 的中点.
\begin{enumerate}
    \item 求证： $ {CE}// $ 平面 $ {PAB} $ ；
    \item 求证： $ l//{AD} $ ；
    \item 若 $ M $ 是线段 $ {CE} $ 上一动点，则线段 $ {AD} $ 上是否存在点 $ N $ ， 使 $ {MN}// $ 平面 $ {PAB} $ ? 说明理由.   
\end{enumerate}

\begin{center}

    \includegraphics[width=0.25\linewidth]{bodyxb//src/a'osuhf'aiuhf'asf.png}
\end{center}

\begin{jieda}
    \begin{enumerate}
        \item 在四棱锥 $ P - {ABCD} $ 中,取 $ {PA} $ 中点 $ F $ ,连 $ {BF} $ 、 $ {EF} $
$ \because {EF} = \dfrac{1}{2}{AD},{EF} ⫫ {AD} $
又 $ \because {BC} = \dfrac{1}{2}{AD},{BC} ⫫ {AD} $
$ \therefore {EF} = {BC},{EF}\parallel {BC} $
$ \therefore $ 四边形 $ {BCEF} $ 为平行四边形,
$ \therefore {CE} ⫫ {BF} $ ,
又 $ \because {CE} ⊄ $ 平面 $ {PAB},{BF} \subset $ 平面 $ {PAB},\therefore {CE} \cup $ 平面 $ {PAB} $ ;

\begin{center}

    \includegraphics[width=0.25\linewidth]{bodyxb//src/asufgh'aiuhf'aiuhsfd.png}
\end{center}

        \item $ \because $ 在梯形 $ {ABCD} $ 中, $ {BC}\square {AD} $ ,
又 $ {BC} ⊄ $ 面 $ {PAD},{AD} \subset $ 面 $ {PAD} $ ,
$ \therefore {BC} \parallel $ 面 $ {PAD} $ ,
$ \because {BC} \subset $ 面 $ {PBC} $ ,面 $ {PBC} \cap $ 面 $ {PAD} = l $
$ \therefore {BC}\parallel l $
$ \because {BC}\parallel {AD}\;\therefore {l}^{\prime \prime }{AD} $
        
        \item 取 $ {AD} $ 中点 $ N $ ,连接 $ {CN} $ , $ {EN} $ ,
$ \because E,N $ 分别为 $ {PD},{AD} $ 的中点,
$ \therefore {EN}\parallel {PA} $ ,
$ \because {EN} ⊄ $ 平面 $ {PAB},{PA} \subset $ 平面 $ {PAB} $ ,
$ \therefore {EN} \parallel $ 平面 $ {PAB} $ ,
又由 (1) 可得 $ {CE}\parallel $ 平面 $ {PAB},{CE} \cap  {EN} = E,{CE}\text{、}{EN} \subset $ 平面 $ {CEN} $
$ \therefore $ 平面 $ {CEN} ⫫ $ 平面 $ {PAB} $ ,
$ \because M $ 是 $ {CE} $ 上的动点, $ {MN} \subset $ 平面 $ {CEN} $ ,
$ \therefore {MN} ⫫ $ 平面 $ {PAB} $ ,
$ \therefore $ 当 $ N $ 为中点时, $ {MN} \parallel $ 平面 $ {PAB} $ .
    \end{enumerate}
\end{jieda}

\item 如图1，等腰梯形 $ {ABCD} $ 中， $ {AD}//{BC} $ ， $ E $ 为 $ {AD} $ 边上一点，且 $ {AB} = {AE} = {BE} = {CD} = 2 $ ， $ {BC} = {ED} = 4,O $ 为 $ {BE} $ 中点， $ F $ 为 $ {BC} $ 中点. 将 $ \bigtriangleup  {ABE} $ 沿 $ {BE} $ 折起到 $ \bigtriangleup  {{A}^{\prime }{BE}} $ 的位置，如图 2 .

\begin{center}

    \includegraphics[width=0.25\linewidth]{bodyxb//src/aiushf'aih'fudafsd.png}
\end{center}

\begin{enumerate}
    \item 证明： $ {CD} \bot $ 平面 $ {A}^{\prime }{OF} $ ；
    \item 若 $ {A}^{\prime }O \bot  {OF} $ ，求点 $ F $ 到平面 $ {A}^{\prime }{EC} $ 的距离；
    \item 试在线段 $ {A}^{\prime }D $ 上确定一点 $ G $ ,使得平面 $ {GCE}// $ 平面 $ {A}^{\prime }{OF} $ ,并给出证明.
\end{enumerate}
    \begin{jieda}
        \begin{enumerate}
            \item 在等腰梯形 $ {ABCD} $ 中, $ {AD}//{BC} $ , $ {BC} = {ED} = 4 $ ,
则四边形 $ {BCDE} $ 是平行四边形,则 $ {BE}//{CD} $ ,
因为 $ {AB} = {AE} = {BE} = {CD} = 2 $ ,所以 $ \Delta {ABE} $ 为等边三角形,则 $ \angle {BAE} = {60}^{ \circ  } $
因为 $ O $ 为 $ {BE} $ 中点,所以 $ {BE}\bot {AO} $ ,
在等腰梯形 $ {ABCD} $ 中,可得 $ \angle {BAE} = \angle {CDE} = \angle {CBE} = {60}^{ \circ  } $
连接 $ {EC} $ ,在 $ {\Delta BEC} $ 中,由余弦定理可 $ E{C}^{2} = B{C}^{2} + B{E}^{2} - {2BC} \cdot  {BE} \cdot  \cos \angle {EBC} = {12} $ ,
则 $ {EC} = 2\sqrt{3} $ ,所以 $ B{E}^{2} + E{C}^{2} = {16} = B{C}^{2} $ ,则 $ {BE}\bot {EC} $ .
因为 $ O $ 、 $ F $ 分别是 $ {BE} $ 、 $ {BC} $ 中点,所以 $ {OF}//{EC} $ ,所以 $ {OF}\bot {BE} $ ,
从而可得 $ {OF} \bot  {CD},{CD} \bot  {AO} $ ,
因为 $ \dot{AO} \cap  {OF} = O,\dot{AO}\text{、}{OF} \subset $ 平面 $ \dot{AOF} $ ,所以 $ {CD} \bot $ 平面 $ \dot{AOF} $ ；
    \item 由(1)可知 $ {OF}//{EC} $ ,因 $ {OF} ⊄ $ 平面 $ {AEC} $ , $ {EC} \subset $ 平面 $ {AEC} $ ,则 $ {OF}// $ 平面 $ {AEC} $ ,
所以点 $ F $ 到平面 $ A\dot{}{EC} $ 的距离即为点 $ O $ 到平面 $ A\dot{}{EC} $ 的距离,
因为 $ O $ 是 $ {BE} $ 中点,所以点 $ F $ 到平面 $ {AEC} $ 的距离等于点 $ B $ 到平面 $ {AEC} $ 的距离的一半.
取 $ {A}^{\prime }E $ 的中点为 $ H $ ,连接 $ {BH} $ ,
因为 $ \Delta {A}^{\prime }{BE} $ 为等边三角形,所以 $ {BH} \bot  {A}^{\prime }E $ ,

\begin{center}

    \includegraphics[width=0.25\linewidth]{bodyxb//src/asfuha'iuhf'aushdfa.png}
\end{center}


由(1)知 $ {EC}\bot {BE} $ ,因为 $ {A}^{\prime }O\bot {OF} $ ,而 $ {OF}//{EC} $ ,则 $ {EC}\bot {A}^{\prime }O $ ,
又因为 $ {BE} \cap  {AD} = O,{BE}\text{、}{AD} \subset $ 平面 $ {ABE} $ ,所以 $ {EC}\bot $ 平面 $ {ABE} $ ,
因为 $ {BH} \subset $ 平面 $ {A}^{\prime }{BE} $ ,所以 $ {EC} \bot  {BH} $ ,
又因为 $ {EC} \cap  {AE} = E,{EC}\text{、}{AE} \subset $ 平面 $ {AEC} $ ,所以 $ {BH}\bot $ 平面 $ {AEC} $ ,
因为 $ \Delta {A}^{\prime }{BE} $ 是等边三角形,边长为 2,故 $ {BH} = \sqrt{{2}^{2} - {1}^{2}} = \sqrt{3} $ ,
所以点 $ B $ 到平面 $ {A}^{\prime }{EC} $ 的距离为 $ \sqrt{3} $ ,
故点 $ F $ 到平面 $ {AEC} $ 的距离为 $ \dfrac{\sqrt{3}}{2} $ ;
    \item 设 $ {DF} \cap  {CE} = M $ ,则 $ \dfrac{FM}{MD} = \dfrac{CF}{ED} = \dfrac{1}{2} $ ,
在 $ {AD} $ 上取点 $ G $ ,使 $ \dfrac{{A}^{\prime }G}{GD} = \dfrac{1}{2} $ ,从而 $ {GM}//{A}^{\prime }F $ ,
连接 $ {EG} $ ,因为 $ \dot{AF} ⊄ $ 平面 $ {GEC},{GM} \subset $ 平面 $ {GEC} $ ,所以 $ \dot{AF}// $ 平面 $ {GEC} $ ,

\begin{center}

    \includegraphics[width=0.25\linewidth]{bodyxb//src/ufha'shf'aosuhfa'oud.png}
\end{center}
   又由 (1) 可知, $ {OF}//{EC},{OF} ⊄ $ 平面 $ {GEC},{GM} \subset $ 平面 $ {GEC} $ ,
则 $ {OF}// $ 平面 $ {GEC} $ ,
又因为 $ {AF} \cap  {OF} = F,{AF},{OF} \subset $ 平面 $ {AOF} $ ,所以平面 $ {GCE}// $ 平面 $ {AOF} $ ,
故 $ G $ 为线段 $ {AD} $ 上靠近 $ A $ 的三等分点时,平面 $ {GCE}// $ 平面 $ {AOF} $ .

        \end{enumerate}
    \end{jieda}



    
\end{enumerate}



\subsection*{平面与平面的位置关系（2）}


A 组：

\begin{enumerate}


\item 二面角的取值范围是\blkda{$[0,\pi]$}.

\item 二面角是指\dotfill{(\kh{D})}

\xx{两个平面相交所组成的图形}
{一个平面绕这个平面内一条直线旋转所组成的图形}
{从一个平面内的一条直线出发的一个半平面与这个平面所组成的图形}
{从一条直线出发的两个半平面所组成的图形}

\item 如图，已知 $ {PA} \bot  a $ ， $ {PB} \bot  \beta $ ，垂足为 $ {A} $ 、 $ B $ ，若 $ \angle {APB} = {60}^{ \circ  } $ ，则二面角 $ a - l - \beta $ 的大小是\blkda{$\dfrac{2\pi}{3}$}.

\begin{center}

    \includegraphics[width=0.25\linewidth]{bodyxb//src/a'oushfaouhf'shdfas.png}
\end{center}

    
\end{enumerate}

B 组:

\begin{enumerate}


\item 若分别向小于 90° 的二面角的两个半平面作垂线，且它们是异面直线，则这两条直线所成的角与二面角的平面角的关系是 \dotfill{(\kh{A})}

\xx{相等} {互补} {互余} {既不相等又不互补也不互余}



\item 已知三棱锥 $ D - {ABC} $ 的三个侧面与底面全等，且 $ {AB} = {AC} = \sqrt{3},{BC} = 2 $ ，则二面角 $ D - {BC} - A $ 的余弦值为\blkda{0}.


\item 在正三角形 $ {ABC} $ 中， $ {AD}\bot {BC} $ 于 $ D $ ，沿 $ {AD} $ 折成二面角 $ B $ — $ {AD} - C $ 后， $ {BC} = \dfrac{1}{2}{AB} $ ， 这时二面角 B-AD-C 的大小为\blkda{$\dfrac{\pi}{3}$}.


\item 已知 $ P $ 为锐二面角内一点，且 $ P $ 到两个半平面及棱的距离之比为 $ 1 : \sqrt{2} : 2 $ ，则此二面角的度数为\blkda{$75^{\circ}$}.


\item 在 ABCD 为正方形，E 为 AB 中点，将 $ \bigtriangleup  {DAE} $ 和 $ \bigtriangleup  {CBE} $ 分别沿 DE、CE 折起， 使 A 与 B 重合, 重合后的点记为 P, 则平面 PCD 与平面 ECD 所成二面角的大小为\blkda{$\dfrac{\pi}{6}$}.



\item 已知菱形的一个内角是 $ {60}^{ \circ  } $ ，边长为 $ a $ ，沿菱形较短的对角线折成大小为 $ {60}^{ \circ  } $ 的二面角，则菱形中含 $ {60}^{ \circ  } $ 角的两个顶点间的距离为
\blkda{$\dfrac{\sqrt{3}a}{2}$}.


\item 在正方体 $ {{ABCD}} - {{A}}_{1}{{\;B}}_{1}{{C}}_{1}{{D}}_{1} $ 中，平面 $ {{A}}_{1}{{{DC}}}_{1} $ 和平面 $ {{A}}_{1}{{{BC}}}_{1} $ 所成的二面角 (锐角) 大小为\blkda{$arccos\left(\dfrac{1}{3}\right)$}.




\item 已知二面角 $ a - {AB} - \beta $ 的平面角是锐角 $ \theta $ ， $ a $ 内一点C到 $ \beta $ 的距离为3，点C 到棱AB的距离为4，则 $ \tan \theta  = $ \blkda{$\dfrac{3\sqrt{7}}{7}$}.



\item 在 30° 的二面角的一个面内有一个点，若它到另一个面的距离是 10，则它到棱的距离为\blkda{20}.



\item 若二面角内一点到二面角的两个面的距离分别为 $ a $ 和 $ \sqrt{2}a $ ，到棱的距离为 $ {2a} $ ，则此二面角的度数是\blkda{$75^{\circ}$或$165^{\circ}$}.



\item 在棱长为 $a$ 的正方体 $ {{ABCD}} - {{A}}_{1}{{B}}_{1}{{C}}_{1}{{D}}_{1} $ 中，求二面角 $ {{C}}_{1} - {{BD}} - {C} $ 的大小.

\begin{jieda}
    按照定义法，$\angle{C_1DC}$即为二面角的平面角，因此$arctan\sqrt{2}$.
\end{jieda}

\item 过正方形 $ {ABCD} $ 的顶点 $ A $ 作 $ {PA} \bot $ 平面 $ {ABCD} $ ，设 $ {PA} = {AB} = a $ ，求二面角 $ B - PC-D$ 的大小.

\begin{jieda}
    连接 $ {AC},{BD} $ ,因为 $ {PA} \bot $ 平面 $ {ABD},{BD} \subset $ 平面
$ {ABD} $ ,所以 $ {PA} \bot  {BD} $ ,又 $ {BD} \bot  {AC},{AC},{PA} \subset $ 平面 $ {PAC} $ ,
$ {PA} \cap  {AC} = A $ ,所以 $ {BD}\bot $ 平面 $ {PAC} $ ,又 $ {PC} \subset $ 平面 $ {PAC} $ ,
所以 $ {BD} \bot  {PC} $ ,过 $ {BD} $ 作平面 $ {BDH} \bot  {PC} $ 于 $ H $ ,则 $ {PC} \bot $
$ {DH},{PC} \bot  {BH} $ ,所以 $ \angle {BHD} $ 为二面角 $ B - {PC} - D $ 的平
面角. 由题知 $ {PB} = \sqrt{2}a,{BC} = a $ ,所以 $ {PC} = \sqrt{3}a,\dfrac{1}{2}{PB} $
$ \cdot  {BC} = {S}_{\bigtriangleup {PBC}} = \dfrac{1}{2}{PC} \cdot  {BH} $ ,所以 $ {BH} = \dfrac{\sqrt{6}}{3}a = {DH} $ ,又
$ {BD} = \sqrt{2}a $ ,在 $ \bigtriangleup {BHD} $ 中,由余弦定理得 $ \cos \angle {BHD} = $
$ \dfrac{B{H}^{2} + D{H}^{2} - B{D}^{2}}{{2BH} \cdot  {DH}} = \dfrac{{\left( \dfrac{\sqrt{6}}{3}a\right) }^{2} + {\left( \dfrac{\sqrt{6}}{3}a\right) }^{2} - {\left( \sqrt{2}a\right) }^{2}}{2 \times  \dfrac{\sqrt{6}}{3}a \times  \dfrac{\sqrt{6}}{3}a} = $
$ - \dfrac{1}{2} $ ,又 $ 0 < \angle {BHD} < \pi $ ,所以 $ \angle {BHD} = \dfrac{2\pi }{3} $ ,故二面角
$ B - {PC} - D $ 的大小是 $ \dfrac{2\pi }{3} $ .

\end{jieda}



\item 如图，已知 $ \operatorname{Rt} \bigtriangleup  {ABC} $ 的直角边 $ {AB} = 4 $ ， $ {AC} = 2 $ ， $ {PA}\bot $ 平面 $ {ABC} $ ， $ {PA} = 2 $ ，求二面角 $ P - {BC} - A $ 的大小.

\begin{center}

    \includegraphics[width=0.25\linewidth]{bodyxb//src/asuhf'aouhfasfa.png}
\end{center}

\begin{jieda}
如图,过点 $ A $ 作 $ {AD}\bot {BC} $ 于 $ D $ ,连接 $ {PD} $ .

  \begin{center}
  
      \includegraphics[width=0.25\linewidth]{bodyxb//src/hf'ahf'ahf;haosfhasfa.png}
  \end{center}
  $ \because {PA} \bot $ 平面 $ {ABC},\therefore {PA} \bot  {BC} $ ,
又 $ {AD} \bot  {BC},\therefore {BC} \bot $ 平面 $ {PAD},\therefore {BC} \bot  {PD} $ ,
$ \therefore \angle {PDA} $ 是二面角 $ P - {BC} - A $ 的平面角.
$ \because {AB} = 4,{AC} = 2,\therefore {BC} = \sqrt{{4}^{2} + {2}^{2}} = 2\sqrt{5},{AD} = \dfrac{4 \times  2}{2\sqrt{5}} = \dfrac{4\sqrt{5}}{5} $ ,
$ \therefore $ 在 Rt $ \bigtriangleup  {PAD} $ 中, $ \tan \angle {PDA} = \dfrac{PA}{AD} = \dfrac{\sqrt{5}}{2} $ ,
$ \therefore \angle {PDA} = \arctan \dfrac{\sqrt{5}}{2} $ ,即二面角 $ P - {BC} - A $ 的大小为 $ \arctan \dfrac{\sqrt{5}}{2} $ .

\end{jieda}



\item 已知直角三角形 $ {ABC} $ 的斜边 $ {AB} $ 在平面 $ a $ 内， $ {AC} $ 、 $ {BC} $ 与平面 $ a $ 所成的角分别为 $ {30}^{ \circ  } $ 和 $ {45}^{ \circ  } $ ,求 $ \bigtriangleup {ABC} $ 所在平面与 $ a $ 所成的二面角 (锐角) 的大小.

\begin{jieda}
    过点 $ C $ 作 $ {CO}\bot \alpha $ ,交 $ \alpha $ 于 $ O $ ,连结 $ {AO},{BO} $ ,作 $ {CD}\bot l $ ,交 $ l $ 于点 $ D $ ,连结 $ {OD} $ ,
则 $ \angle {CDO} $ 是二面角 $ \alpha  - l - \beta $ 的平面角,
$ \because {{Rt}}\angle {{ABC}} $ 在平面 $ \beta $ 内, $ \angle {ACB} = {90}^{ \circ  } $ ,斜边 $ {AB} $ 在二面角 $ \alpha  - l - \beta $ 的棱 $ l $ 上,
$ {AC} $ 与平面 $ \alpha $ 所成角为45 $ {}^{ \circ  } $ , $ {BC} $ 与平面 $ \alpha $ 所成角为30 $ {}^{ \circ  } $ ,
$ \therefore \angle {CAO} = {45}^{ \circ  },\angle {CBO} = {30}^{ \circ  } $ ,
设 $ {CO} = 1 $ ,则 $ {AO} = 1 $ , $ {AC} = \sqrt{2} $ , $ {BC} = 2 $ , $ {BO} = \sqrt{3} $ , $ {AB} = \sqrt{6} $ ,
$ {CD} = \dfrac{{AC} \times  {BC}}{AB} = \dfrac{\sqrt{2} \times  2}{\sqrt{6}} = \dfrac{2}{\sqrt{3}},{AD} = \sqrt{2 - \dfrac{4}{3}} = \sqrt{\dfrac{2}{3}},{OD} = \sqrt{1 - \dfrac{2}{3}} = \sqrt{\dfrac{1}{3}}, $
$ \therefore \cos \angle {CDO} = \dfrac{C{D}^{2} + O{D}^{2} - C{O}^{2}}{2 \times  {CD} \times  {OD}} = \dfrac{\dfrac{4}{3} + \dfrac{1}{3} - 1}{2 \times  \sqrt{\dfrac{4}{3}} \times  \sqrt{\dfrac{1}{3}}} = \dfrac{1}{2} $ ,
$ \therefore \angle {CDO} = {60}^{ \circ  } $ ,
$ \therefore $ 二面角 $ \alpha  - l - \beta $ 的平面角大小为 $ {60}^{ \circ  } $ .

\begin{center}

    \includegraphics[width=0.25\linewidth]{bodyxb//src/iisgjijgoiaejgofaerfa.png}
\end{center}

\end{jieda}


\item 如图，四棱锥 $ P - {ABCD} $ 的底面是边长为 $ a $ 的菱形， $ \angle {BCD} = {120}^{ \circ  } $ ， $ {PC} \bot $ 平面 $ {ABCD},{PC} = a,E $ 为 $ {PA} $ 的中点， $ O $ 为底面对角线的交点；

\begin{enumerate}
    \item 求证： $ {AC} \bot $ 平面 $ {EDB} $ ；
    \item 求二面角 $ P - {AB} - C $ 的正切值.
\end{enumerate}

\begin{center}

    \includegraphics[width=0.25\linewidth]{bodyxb//src/fa'jshf'ahsfasdf.png}
\end{center}


\begin{jieda}
    \begin{enumerate}
        \item 如图,连接 $ {EO} $ , $ \because E $ 为 $ {PA} $ 的中点, $ O $ 为菱形对角线的交点, $ \therefore O $ 为 $ {AC} $ 中点, $ \therefore {EO} \parallel  {PC} $ ,又 $ \because {PC} \bot $ 平面 $ {ABCD} $
$ \therefore {PC} \bot  {AC} $ , $ \therefore {EO} \bot  {AC} $ ,又 $ \because $ 菱形的对角线互相垂直平分, $ \therefore {AC} \bot  {BD} $ ,又 $ {EO} \cap  {BD} = O $ , $ \therefore {AC} \bot $ 平面 $ {EDB} $ ；

\begin{center}

    \includegraphics[width=0.25\linewidth]{bodyxb//src/dsghjsoijgaoiefaf.png}
\end{center}

        \item 如图,设 $ {AB} $ 中点为 $ F $ ,由底面 $ {ABCD} $ 为菱形可知, $ \bigtriangleup  {ABC} $ 为等边三角形, $ \therefore {CF} \bot  {AB} $ ,
又 $ \because {PC} \bot $ 平面 $ {ABCD} $ , $ \therefore {PC} \bot  {AB} $ , $ \therefore {AB} \bot $ 平面 $ {PCF} $ , $ \therefore {AB} \bot  {PF} $ , $ \therefore \angle {PFC} $ 为二面角 $ P - {AB} - C $ 的平面角, $ {CF} = \dfrac{\sqrt{3}}{2}a - {PC} = a $
,故 $ \tan \angle {PFC} = \dfrac{PC}{CF} = \dfrac{2\sqrt{3}}{3} $.

\begin{center}

    \includegraphics[width=0.25\linewidth]{bodyxb//src/gijsdghsodhgaouhgf.png}
\end{center}

    \end{enumerate}
\end{jieda}

\end{enumerate}



C 组:

\begin{enumerate}

\item 如图，在正方体 $ \mathbf{{ABCD} - {A}_{1}{B}_{1}{C}_{1}{D}_{1}} $ 中，动点 $ \mathbf{M} $ 在线段 $ \mathbf{{A}_{1}C} $ 上，异面直线 $ \mathbf{A}{\mathbf{D}}_{\mathbf{1}} $ 和 $ \mathbf{B}\mathbf{M} $ 所成的角为 $ \mathbf{\theta } $ ，则 $ \mathbf{\theta } $ 的取值范围是\blkda{$[\dfrac{\pi}{6},\dfrac{\pi}{3}]$}.（用区间表示）

\begin{center}

    \includegraphics[width=0.25\linewidth]{bodyxb//src/a'sfh'aouhfadfasfsfa.png}
\end{center}



\item 已知正方体 $ {ABCD} - {A}_{1}{B}_{1}{C}_{1}{D}_{1} $ 的棱长为 $ 4,M $ 为棱 $ {DC} $ 的中点， $ N $ 为侧面 $ B{C}_{1} $ 的中心，过点 $ M $ 的平面 $ \alpha $ 垂直于 $ {DN} $ ，则平面 $ \alpha $ 截正方体 $ A{C}_{1} $ 所得的截面面积为\blkda{$4\sqrt{6}$}.




\end{enumerate}





\subsection*{平面与平面的位置关系（3）}


A 组:

\begin{enumerate}




\item 下列命题正确的是\dotfill{(\kh{B})}

\xx{直线 $ a// $ 平面 $ a $ ,直线 $ b\bot a,b \subset $ 平面 $ \beta $ ,则 $ a \bot  \beta $}
{若直线 $ a \bot $ 直线 $ b,a \bot $ 平面 $ a,b \bot $ 平面 $ \beta $ ，则 $ a \bot  \beta $}
{过平面外的一条直线有且只有一个平面与已知平面垂直}
{过平面外一点有且只有一个平面与已知平面垂直}

\item 若平面 $ \alpha $ 、 $ \beta $ 互相垂直，则 \dotfill{(\kh{D})}

\xx{$ \alpha $ 中的任意一条直线垂直于 $ \beta $} 
{$ \alpha $ 中有且只有一条直线垂直于 $ \beta $}
{平行于 $ \alpha $ 的直线垂直于 $ \beta $} 
{$ \alpha $ 内垂直于交线的直线必垂直于 $ \beta $}


\item 已知三个平面 $ a $ 、 $ \beta $ 、 $ \gamma $ ，则下列命题中正确的个数是 \dotfill{(\kh{A})}
\begin{enumerate}

\item 若 $ a \bot  \beta ,\gamma  \bot  \beta $ ，则 $ a \bot  \gamma $ ；
\item 若 $ a \bot  \beta ,\gamma  \bot  \beta , $ ，则 $ a//\gamma $ ；
\item 若 $ a//\beta ,\beta  \bot  \gamma $ ，则 $ a \bot  \gamma $ ；
    
\end{enumerate}
\xx{1}{2}{3}{0} 

\item 已知两个平面互相垂直，那么下列命题中正确的个数是 \dotfill{(\kh{C})}
\begin{enumerate}

\item 一个平面内的直线必垂直于另一个平面内的无数条直线；

\item 一个平面内且垂直于这两个平面交线的直线必垂直于另一个平面内的任意一条直线;

\item 一个平面内的任何一条直线必垂直于另一个平面；

\item 过一个平面内的任意一点作交线的垂线，则此直线必垂直于另一个平面。
\end{enumerate}
\xx{4}{3}{2}{1}


\item 在空间两两垂直的平面最多有\dotfill{(\kh{B})}

\xx{两个} {三个} {四个} {无数个}



\end{enumerate}


B 组:

\begin{enumerate}


\item 在空间四边形 ABCD 中，已知二面角 A一BCD,A一CD=B,A一BD一C 都相等，则 A 点在平面 BCD 上的射影是 $ \bigtriangleup  {BCD} $ 的\dotfill{(\kh{A})}

\xx{内心}{外心} {垂心} {重心}



\item 若一个二面角的两个半平面分别垂直与另一个二面角的两个半平面，则这两个二面角的关系是\dotfill{(\kh{D})}

\xx{相等} {互补} {相等或互补} {不能确定}

\item 从二面角 P-MN—Q 内一点 A 分别作 AB⊥平面 P，AC⊥平面 Q，其中 B、C 为垂足， $ \angle {BAC} = {60}^{ \circ  } $ ，那么二面角 P-MN—Q 的大小是\blkda{$\dfrac{2\pi}{3}$}.

\item 已知 $ {A} $ 、 $ {B} $ 是二面角 $ a - l - \beta $ 的棱 $ l $ 上两点，P 是面 $ \beta $ 内一点，PB $ \bot  l $ 于点 B， PA 和 $ l $ 所成的角为 $ {45}^{ \circ  } $ ，PA 和面 $ a $ 所成的角为 $ {30}^{ \circ  } $ ，则二面角 $ a - l - \beta $ 的大小为\blkda{$\dfrac{\pi}{4}$或$\dfrac{3\pi}{4}$}.

\item 已知点 $ {A} $ 、 $ {B} $ 分别是二面角 $ {A} - {{CD}} - {B} $ 的各面上的一点， $ {{AC}} \bot  {{CD}} $ ， $ {{BD}} \bot  {{CD}} $ ，C、D 分别为垂足，若 CD=3,AC= $ \sqrt{2} $ ,BD=2,AB= $ \sqrt{10} $ ，则二面角 A—CD—B 的度数是\blkda{$arccos\dfrac{5\sqrt{2}}{8}$}.



\item 如图，矩形 ABEF 和正方形 ABCD 有公共边 AB，它们所在的半平面成 $ {60}^{ \circ  } $ 的二面角,且 $ {{AB}} = {2a},{{BF}} = a $ ,则 $ {{FC}} = $\blkda{$\sqrt{7}a$}


\begin{center}

    \includegraphics[width=0.25\linewidth]{bodyxb//src/f'aohf'ajhfasfa.png}
\end{center}



\item 在正方体 $ {{ABCD}} - {{A}}_{1}{{B}}_{1}{{C}}_{1}{{D}}_{1} $ 中，平面 $ {{{AA}}}_{1}{{C}}_{1}{C} $ 和平面 $ {{A}}_{1}{{{BCD}}}_{1} $ 所成的二面角（锐角）大小为\blkda{$\dfrac{\pi}{3}$}.


\item 在矩形 $ {ABCD} $ 中， $ {AB} = 3,{AD} = 4,{PA} \bot $ 平面 $ {ABCD},{PA} = \sqrt{3} $ ，则平面 $ {PBD} $ 与平面 $ {ABCD} $ 的夹角的正切值为\blkda{$\dfrac{5\sqrt{3}}{12}$}.


\item 如图，在空间四边形 $ {ABCD} $ 中， $ \bigtriangleup  {BCD} $ 是正三角形， $ \bigtriangleup {ABD} $ 是等腰直角三角形，且 $ \angle {BAD} = {90}^{ \circ  } $ ，又二面角 $ A - {BD} - C $ 为直二面角,则二面角 $ A - {CD} - B $ 的正切值为\blkda{$\dfrac{2\sqrt{3}}{3}$}.
\begin{center}

    \includegraphics[width=0.25\linewidth]{bodyxb//src/hf'ihfuu'shdfui'sdfh.png}
\end{center}



\item 在二面角 $ \alpha  - {AB} - \beta $ 的一个平面 $ \alpha $ 内，有一直线 $ {AC} $ ，它与棱 $ {AB} $ 成 $ {45}^{ \circ  } $ 角， $ {AC} $ 与平面 $ \beta $ 成 $ {30}^{ \circ  } $ 角,二面角 $ \alpha  - {AB} - \beta $ 的大小为\blkda{$\dfrac{\pi}{4}$或$\dfrac{3\pi}{4}$}





\item 已知 $ {P} $ 为 $ {120}^{ \circ  } $ 的二面角 $ a - a - \beta $ 内一点，且点 $ {P} $ 到平面 $ a $ 和 $ \beta $ 的距离分别为 6 和8，求点 $ {P} $ 到棱 $ a $ 的距离.
\begin{jieda}
    设长度为x，过P作$\alpha$和$\beta$的垂线，交于A、B两点，过P作棱的垂线，交于M，则有$sin\angle{PAM}=\dfrac{6}{x}$,$sin\angle{PBM}=sin(120^{\circ}-\angle{PAM})=\dfrac{8}{x}$,从而求得$x=\dfrac{4\sqrt{39}}{3}$.
\end{jieda}



\item 如图，在三棱锥 $ P - {ABC} $ 中， $ {AB}\bot {BC} $ ，

$ {AB} = {PB} = {PC} = 2,{PA} = {BC} = 2\sqrt{2} $ ，棱 $ {PB},{PA},{AC},{BC} $ 的中点分别为 $ D,E,F,O $ .

\begin{enumerate}
    \item 证明： $ {EF}// $ 平面 $ {ADO} $ ；
    \item 证明：平面 $ {BEF} \bot $ 平面 $ {PAO} $ .
\end{enumerate}

\begin{center}

    \includegraphics[width=0.25\linewidth]{bodyxb//src/foushf'oushf'ouhsdfsdfs.png}
\end{center}

\begin{jieda}
    \begin{enumerate}
        \item 因为棱 $ {PB},{PA},{AC},{BC} $ 的中点分别为 $ D,E,F,O $ ,
所以 $ {EF}//{PC},{DO}//{PC} $ ,所以 $ {EF}//{DO} $ ,
又因为 $ {DO} \subset $ 平面 $ {ADO} $ , $ {EF} ⊄ $ 平面 $ {ADO} $ ,所以 $ {EF}// $ 平面 $ {ADO} $ .
        \item 因为 $ {PB} = {PC} = 2,{BC} = 2\sqrt{2} $ , $ O $ 是 $ {BC} $ 的中点,
所以 $ {PO}\bot {BC} $ ,且 $ {PO} = \sqrt{2} $ .
因为 $ {AB}\bot {BC} $ ,所以 $ {AO} = \sqrt{A{B}^{2} + B{O}^{2}} = \sqrt{6} $ ,
所以 $ P{O}^{2} + A{O}^{2} = P{A}^{2} $ ,所以 $ {PO}\bot {AO} $ .
又因为 $ {AO} \cap  {BC} = O,{AO},{BC} \subset $ 平面 $ {ABC} $ ,所以 $ {PO}\bot $ 平面 $ {ABC} $ .
又 $ {BF} \subset $ 平面 $ {ABC} $ ,所以 $ {PO} \bot  {BF} $ .
如图,连接 $ {OF} $ ,则 $ {OF} \parallel  {AB} $ ,从而 $ {OF}\bot {BC} $ ,且 $ {OF} = \dfrac{1}{2}{AB} = 1 $ ,所以 $ \dfrac{OB}{AB} = \dfrac{OF}{OB} = \dfrac{\sqrt{2}}{2} $ ,得 $ \operatorname{Rt} \bigtriangleup  {ABO} \backsim  \operatorname{Rt} \bigtriangleup  {BOF} $ ,
所以 $ \angle {FBO} + \angle {AOB} = {90}^{ \circ  } $ ,所以 $ {AO}\bot {BF} $ .
由 $ {PO}\bot {BF},{AO}\bot {BF},{AO} \cap  {PO} = O,{AO},{PO} \subset $ 平面 $ {POA} $ ,得 $ {BF}\bot $ 平面 $ {POA} $ ,
又 $ {BF} \subset $ 平面 $ {BEF} $ ,所以平面 $ {BEF} \bot $ 平面 $ {POA} $ .

\begin{center}

    \includegraphics[width=0.25\linewidth]{bodyxb//src/fhs'ofhaoighaouhfaohfa.png}
\end{center}

    \end{enumerate}
\end{jieda}




\item 如图，在四棱锥 $ M - {ABCD} $ 中，底面 $ {ABCD} $ 是边长为 2 的菱形， $ {MA} \bot $ 平面 $ {ABCD} $ ，平面 $ {MAB} \cap $ 平面 $ {MCD} = l $ ，平面 $ {MAB} \bot $ 平面 $ {MBC} $ .
\begin{enumerate}

\item 证明： $ l //{CD} $ ；

\item 证明： $ {AD} \bot $ 平面 $ {MAB} $ ；

\item 当 $ {MA} $ 为何值时，二面角 $ B - {MC} - D $ 的余弦值为 $ - \dfrac{1}{3} $ .
    
\end{enumerate}

\begin{center}

    \includegraphics[width=0.25\linewidth]{bodyxb//src/glahgd'ush'fouhsdgfa.png}
\end{center}

\begin{jieda}
    \begin{enumerate}
        \item 菱形 $ {ABCD} $ 中, $ {CD}//{AB},{AB} \subset $ 平面 $ {MAB},{CD} ⊄ $ 平面 $ {MAB} $ ,
则 $ {CD}// $ 平面 $ {MAB} $ ,
而平面 $ {MAB} \cap $ 平面 $ {MCD} = l,{CD} \subset $ 平面 $ {MCD} $ ,
所以 $ l//{CD} $ .

        \item 在平面 $ {MAB} $ 内过点 $ {A作AE}\bot {MB} $ 于 $ E $ ,平面 $ {MAB}\bot $ 平面 $ {MBC} $ ,平面 $ {MAB} \cap $ 平面 $ {MBC} = {MB} $ ,
则 $ {AE} \bot $ 平面 $ {MBC} $ ,而 $ {BC} \subset $ 平面 $ {MBC} $ ,于是 $ {AE} \bot  {BC} $ ,
又 $ {MA}\bot $ 平面 $ {ABCD},{BC} \subset $ 平面 $ {ABCD} $ ,
则 $ {MA}\bot {BC} $ ,而 $ {MA} \cap  {AE} = A,{MA} \subset $ 平面 $ {MAB} $ ,
因此 $ {BC}\bot $ 平面 $ {MAB} $ ,又 $ {BC}//{AD} $ ,
所以 $ {AD} \bot $ 平面 $ {MAB} $ .

\begin{center}

    \includegraphics[width=0.25\linewidth]{bodyxb//src/gsrhg'uahgifa'fhaw.png}
\end{center}

        \item 由(2)知 $ {AD} \bot $ 平面 $ {MAB} $ , $ {AB} \subset $ 平面 $ {MAB} $ ,则 $ {AB} \bot  {AD} $ ,菱形 $ {ABCD} $ 为正方形,
由 $ {BC} \bot $ 平面 $ {MAB},{MB} \subset $ 平面 $ {MAB} $ ,得 $ {BC} \bot  {MB} $ ,
过 $ B $ 作 $ {BO}\bot {MC} $ 于 $ O $ ,连接 $ {BD},{DO} $ ,
$ {MB} = \sqrt{A{B}^{2} + M{A}^{2}} = \sqrt{A{D}^{2} + M{A}^{2}} = {MD} $ ,而 $ {BC} = {DC},{MC} = {MC} $ ,
则 $ {\Delta MBC} \cong  {\Delta MDC} $ ,有 $ \angle {BMO} = \angle {DMO} $ ,于是 $ {\Delta BMO} \cong  {\Delta DMO} $ ,则 $ \angle {DOM} = \angle {BOM} = {90}^{ \circ  } $ ,
即 $ {DO}\bot {MC},\angle {BOD} $ 是二面角 $ B - {MC} - D $ 的平面角,
令 $ {MA} = t,{MC} = \sqrt{M{B}^{2} + B{C}^{2}} = \sqrt{M{A}^{2} + A{B}^{2} + B{C}^{2}} = \sqrt{{t}^{2} + 8} $ ,
$ {DO} = {BO} = \dfrac{{MB} \cdot  {BC}}{MC} = \dfrac{2\sqrt{{t}^{2} + 4}}{\sqrt{{t}^{2} + 8}} $ ,而 $ {BD} = 2\sqrt{2} $ ,在 $ {\Delta BOD} $ 中,由余弦定理得：
$ \cos \angle {BOD} = \dfrac{2{\left( \dfrac{2\sqrt{{t}^{2} + 4}}{\sqrt{{t}^{2} + 8}}\right) }^{2}}{2{\left( \dfrac{2\sqrt{{t}^{2} + 4}}{\sqrt{{t}^{2} + 8}}\right) }^{2}} =  - \dfrac{4}{{t}^{2} + 4} =  - \dfrac{1}{3} $ ,解得 $ t = 2\sqrt{2} $ ,
所以当 $ {MA} $ 的值为 $ 2\sqrt{2} $ 时,二面角 $ B - {MC} - D $ 的余弦值为 $ - \dfrac{1}{3} $ .

    \end{enumerate}
\end{jieda}

    
\end{enumerate}

C 组:
\begin{enumerate}

\item 在直三棱柱 $ {ABC} - {A}_{1}{B}_{1}{C}_{1} $ 中， $ {A{A}_{1}} = 1,\;{{A}_{1}C} = 2,\;{\angle {BAC}} = \dfrac{\pi }{6},\;E,F $ 为线段 $ {A}_{1}C $ 的三等分点,点 $ D $ 在线段 $ {EF} $ 上 (包括端点) 运动,则二面角 $ D - {AB} - C $ 的正弦值的取值范围为\blkda{$[\dfrac{1}{2},\dfrac{4\sqrt{19}}{19}]$}.

\item 已知在矩形 $ {ABCD} $ 中， $ {AB} = 2 $ ， $ {BC} = \sqrt{2} $ ， $ P $ 为 $ {AB} $ 的中点，将 $ \bigtriangleup  {ADP} $ 沿 $ {DP} $ 翻折，得到四棱锥 $ {A}_{1} - {BCDP} $ ，则二面角 $ {A}_{1} - {DC} - B $ 的余弦值最小是\blkda{$\dfrac{\sqrt{2}}{2}$}.



\end{enumerate}


\clearpage

\subsection*{平面与平面的位置关系（4）}

\begin{enumerate}

\item 已知直线 $ l $ ， $ m $ 和平面 $ a $ ， $ \beta $ ，且 $ l \bot  a $ ， $ m \subset  \beta $ ，则下列命题中正确的是\blkda{(1)(4)}. ①若 $ a//\beta $ ，则 $ l \bot  m $ ②若 $ l \bot  m $ ，则 $ a//\beta $ ③若 $ a \bot  \beta $ ，则 $ l//m $ ④若 $ l//m $ ，则 $ a \bot  \beta $

\item 在长方体 $ {ABCD} - {A}_{1}{B}_{1}{C}_{1}{D}_{1} $ 中，过直线 $ A{C}_{1} $ 的平面交直线 $ B{B}_{1} $ 于点 $ E $ ，交直线 $ D{D}_{1} $ 于点 $ F $ ，则四边形 $ {AE}{C}_{1}F $ 的形状为\blkda{平行四边形}.

\item 已知三棱锥 $ D - {ABC} $ 的三个侧面与底面全等，且 $ {AB} = {AC} = \sqrt{3},{BC} = 2 $ ，则二面角 $ D - {BC} - A $ 的余弦值为\blkda{$0$}.


\item 如图，在长方体 $ {ABCD} - {A}_{1}{B}_{1}{C}_{1}{D}_{1} $ 中， $ {AB} = {BC} = \dfrac{\sqrt{6}}{3}{C{C}_{1}} $ ，则二面角 $ {C}_{1} - {DB} - C $ 的大小为\blkda{$\dfrac{\pi}{3}$}.

\begin{center}

    \includegraphics[width=0.25\linewidth]{bodyxb//src/ogh'sdofhg'sdgsda.png}
\end{center}

\item 直二面角 $ a - l - \beta $ 的棱 $ l $ 上有一点 $ A $ ，在平面 $ a,\beta $ 内各有一条射线 $ {AB},{AC} $ 与 $ l $ 均成 $ {45}^{ \circ  } $ ，则 $ \angle {BAC} = $ \blkda{$\dfrac{\pi}{3}$或$\dfrac{2\pi}{3}$}.

\item 若等边三角形 $ {{ABC}} $ 的边长为 1，BC 上的高为 AD，沿高 AD 折成直二面角， 则 A 到 BC 的距离是\blkda{$\dfrac{\sqrt{14}}{4}$}.


\item 已知 $ \bigtriangleup  {{ABC}} $ 和平面 $ \alpha ,\angle {A} = {30}^{ \circ  },\angle {B} = {60}^{ \circ  },{{AB}} = 2,{{AB}} \subset  \alpha $ ， 且平面 $ {ABC} $ 与 $ \alpha $ 所成角为 $ {30}^{ \circ  } $ ，则点 $ C $ 到平面 $ \alpha $ 的距离为\blkda{$\dfrac{\sqrt{3}}{4}$}.


\item 河堤斜面与水平面所成的二面角为 $ {60}^{ \circ  } $ ，堤面上有一条直道 CD，它与坡脚的水平线的夹角为 $ {30}^{ \circ  } $ ，沿这条直道从坡脚向上行走 $ {10}\sqrt{3}m $ 时人升高了\blkda{$\dfrac{15}{2}$m}

\item 已知 $ a $ ， $ \beta $ 是不重合的平面， $ m $ ， $ n $ 是不重合的直线，下列命题中正确的是\dotfill{(\kh{D})}

\xx{若 $ m//a,m//\beta $ ，则 $ a//\beta $}
{若 $ m//n,m//a $ ，则 $ n//a $}
{若 $ m \bot  a,m \bot  \beta $ ，则 $ a \bot  \beta $} 
{若 $ m \bot  a,m//\beta $ ，则 $ a \bot  \beta $}


\item 已知 $ l $ 为空间的一条直线， $ a $ ， $ \beta $ 为空间的两个不同平面，则下列命题错误的是
\dotfill{(\kh{A})}

\xx{若 $ l//a,l \subset  \beta $ ，则 $ a//\beta $}
{若 $ l \bot  a,l \subset  \beta $ ,则 $ a \bot  \beta $}
{若 $ a//\beta ,l \subset  a $ ,则 $ l//\beta $}
{若 $ a \bot  \beta ,a \cap  \beta  = b,l \bot  b,l \subset  a $ ,则 $ l \bot  \beta $}



\item 如图，在四棱锥 $ P - {ABCD} $ 中，底面 $ {ABCD} $ 是直角梯形, $ {AD}//{BC},{AB} \bot  {BC},{AB} = {BC} = 1,{AD} = 2 $ , $ {PA} \bot $ 平面 $ {ABCD},{PA} = 2 $ .

\begin{enumerate}
    \item 求证： $ {CD}\bot $ 平面 $ {PAC} $ ；
    \item 求二面角 $ P - {CD} - A $ 的正弦值.

\end{enumerate}

\begin{center}

    \includegraphics[width=0.25\linewidth]{bodyxb//src/fahsfhasfasfasg.png}
\end{center}

\begin{jieda}
    
\begin{enumerate}
    \item 因为 $ {AB} = {BC} = 1 $ , $ {AB}\bot {BC} $ ,所以 $ {AC} = \sqrt{2} $ , $ \angle {BAC} = {45}^{ \circ  } $ .
又 $ {AB}\bot {BC},{AD} \parallel  {BC} $ ,所以 $ \angle {DAC} = {45}^{ \circ  } $ .
所以 $ C{D}^{2} = C{A}^{2} + D{A}^{2} - {2CA} \cdot  {DA} \cdot  \cos \angle {DAC} = 2 + 4 - 2 \times  \sqrt{2} \times  2 \times  \dfrac{\sqrt{2}}{2} = 2 $ .
所以 $ {CD} = \sqrt{2} $ .
因为 $ {2}^{2} = {\left( \sqrt{2}\right) }^{2} + {\left( \sqrt{2}\right) }^{2} $ ,即 $ A{D}^{2} = A{C}^{2} + C{D}^{2} $ ,
所以 $ \bigtriangleup  {ACD} $ 为直角三角形,且 $ {CD} \bot  {CA} $ .
又 $ {PA}\bot $ 平面 $ {ABCD},{CD} \subset $ 平面 $ {ABCD} $ ,所以 $ {PA}\bot {CD} $ .
$ {PA},{CA} \subset $ 平面 $ {PAC},{PA} \cap  {CA} = A $ ,所以 $ {CD} \bot $ 平面 $ {PAC} $ .
    \item 因为 $ {CD}\bot $ 平面 $ {PAC},{CA},{CP} \subset $ 平面 $ {PAC} $ ,
所以 $ {CD}\bot {CP},{CD}\bot {CA} $ .
所以 $ \angle {PCA} $ 即为二面角 $ P - {CD} - A $ 的平面角.
在Rt $ {\Delta PAC} $ 中, $ {PA}\bot {AC},{PA} = 2,{AC} = \sqrt{2} $ ,所以 $ {PC} = \sqrt{{PA}^{2} + {AC}^{2}} = \sqrt{6} $ ,
所以 $ \sin \angle {PCA} = \dfrac{PA}{PC} = \dfrac{2}{\sqrt{6}} = \dfrac{\sqrt{6}}{3} $ .
即二面角 $ P - {CD} - A $ 的正弦值为 $ \dfrac{\sqrt{6}}{3} $ .


\end{enumerate}
\end{jieda}

\item 过空间一点P的三条直线 PA、PB、PC 不共面，且两两成 $ {60}^{ \circ  } $ 角，求二面角 $ {A} - {{PB}} - {C} $ 的大小.

\begin{jieda}
    根据题意, $ \angle {APB} = \angle {BPC} = \angle {CPA} = {60}^{ \circ  }, $ ,
设 $ {PA} = {PB} = {PC} = a $ ,利用余弦定理计算 $ {AB} $ 、
$ {BC}\text{、}{AC} $ 的长度,得到 $ {AB} = {BC} = {AC} = a $ 。
取 $ {BP} $ 中点 $ E $ ,因为 $ \bigtriangleup  {ABP} $ , 和 $ \bigtriangleup  {BCP} $ , 是等边三角形
,所以 $ {AE} = {CE} = \dfrac{\sqrt{3}}{2}a, $ 。
在 $ \bigtriangleup  {AEC}, $ 中,根据余弦定理计算 $ \cos \angle {AEC}, $ 的值,
得到 $ \cos \angle {AEC} = \dfrac{1}{3} $ 。
答：二面角 $ A - {PB} - C $ 的大小为 $ \arccos \dfrac{1}{3}, $ 。
\end{jieda}



\item 已知直三棱柱 $ {ABC} - {A}_{1}{B}_{1}{C}_{1} $ (侧棱和地面垂直，上下底面是全等的三角形), 所有棱长均为 2,求二面角 $ A - {BC} - {A}_{1} $ 的余弦值大小.

\begin{jieda}
    解：取BC中点O,连结AO、
$ {\mathrm{A}}_{1}\mathrm{O} $ ,
$ \because $ 直三棱柱ABC- $ {\mathrm{A}}_{1}{\mathrm{B}}_{1}{\mathrm{C}}_{1} $ ,所有
棱长均为 2,
$ \therefore \mathrm{{AO}}\bot \mathrm{{BC}},{\mathrm{A}}_{1}\mathrm{O}\bot \mathrm{{BC}} $ ,
$ \therefore \angle {A}_{1}{OA} $ 是二面角 $ A - {BC} - {A}_{1} $ 的
平面角,
$ {AO} = \sqrt{4 - 1} = \sqrt{3},{A}_{1}O = \sqrt{3 + 4} = $
$ \sqrt{7} $ ,
$ \therefore \cos \angle {A}_{1}{OA} = \dfrac{AO}{{A}_{1}O} = \dfrac{\sqrt{3}}{\sqrt{7}} = \dfrac{\sqrt{21}}{7} $ .
$ \therefore $ 二面角 $ \mathrm{A} - \mathrm{{BC}} - {\mathrm{A}}_{1} $ 的余弦值为$\dfrac{\sqrt{21}}{7}$.

\begin{figure}[h]
    \centering
    \includegraphics[width=0.25\linewidth]{bodyxb//src/hoijupijuoih'iugiyguftoyrdst.png}
\end{figure}

\end{jieda}

\item 如图，在四棱锥 $ P - {ABCD} $ 中，底面 $ {ABCD} $ 是正方形，侧面 $ {PAD} \bot $ 底面 $ {ABCD} $ ， 且 $ {PA} = {PD} = \dfrac{\sqrt{2}}{2}{AD} $ ，若 $ E $ 、 $ F $ 分别为 $ {PC} $ 、 $ {BD} $ 的中点.

\begin{enumerate}
    \item 求证： $ {EF}// $ 平面 $ {PAD} $ ；
    \item 求证：平面 $ {ABP}\bot $ 平面 $ {CDP} $ .
\end{enumerate}

\begin{center}

    \includegraphics[width=0.25\linewidth]{bodyxb//src/gijsdohg;dgas.png}
\end{center}
\begin{jieda}
    \begin{enumerate}
        \item 证明：连接AC,
        
        \begin{figure}[h]
            \centering
            \includegraphics[width=0.25\linewidth]{bodyxb//src/sdigjh'aouhrfaoiehFAI.png}
        \end{figure}

        
        因为四边形 $ {ABCD} $ 为正方形,且 $ F $ 为 $ {BD} $ 的中点,所以 $ F $ 为 $ {AC} $ 的中点,
又因为 $ E $ 为 $ {PC} $ 的中点,则 $ {EF}//{PA} $ ,
$ \because {EF} ⊄ $ 平面 $ {PAD},{PA} \subset $ 平面 $ {PAD},\therefore {EF}// $ 平面 $ {PAD} $ .
        \item 证明：因为四边形 $ {ABCD} $ 为正方形,则 $ {CD}\bot {AD} $ ,
因为平面 $ {PAD} \bot $ 平面 $ {ABCD} $ ,平面 $ {PAD} \cap $ 平面 $ {ABCD} = {AD},{CD} \subset $ 平面 $ {ABCD} $ ,
$ \therefore {CD}\bot $ 平面 $ {PAD} $ ,
$ \because {PA} \subset $ 平面 $ {PAD},\therefore {PA} \bot  {CD} $ ,
$ \because {PA} = {PD} = \dfrac{\sqrt{2}}{2}A{D}^{\prime } $ 所以 $ P{A}^{2} + P{D}^{2} = A{D}^{2} $ ,则 $ {PA} \bot  {PD} $ ,
$ \because {CD} \cap  {PD} = D,{CD} \cap  {PD} \subset $ 平面 $ {CDP} $ ,
$ \therefore {PA}\bot $ 平面 $ {CDP} $ ,又 $ {PA} \subset $ 平面 $ {ABP} $ ,
$ \therefore $ 平面 $ {ABP} \bot $ 平面 $ {CDP} $ .
    \end{enumerate}
\end{jieda}



\item 如图所示，在四棱锥 $ P - {ABCD} $ 中，底面 $ {ABCD} $ 是 $ \angle {DAB} = {60}^{ \circ  } $ 且边长为 $ a $ 的菱形，侧面 $ {PAD} $ 为正三角形，其所在的平面垂直于底面 $ {ABCD} $ . 
\begin{enumerate}
    \item 若 $ G $ 为 $ {AD} $ 边的中点，求证： $ {BG}\bot $ 平面 $ {PAD} $ ； 
    \item 求证： $ {AD}\bot {PB} $ ;
    \item 求二面角 $ A - {BC} - P $ 的大小；\item 若 $ E $ 为 $ {BC} $ 边的中点，能否在棱 $ {PC} $ 上找一点 $ F $ ,使得平面 $ {DEF}\bot $ 平面 $ {ABCD} $ ? 并证明你的结论.
\end{enumerate}

\begin{center}

    \includegraphics[width=0.25\linewidth]{bodyxb//src/'doigjsdigjh'gadg.png}
\end{center}

\begin{jieda}
    \begin{enumerate}
        \item 在底面菱形 $ {ABCD} $ 中, $ \angle {DAB} = {60}^{ \circ  } $ , $ G $ 为 $ {AD} $ 边的中点,所以 $ {BG}\bot {AD} $ ,
又平面 $ {PAD} \bot $ 平面 $ {ABCD} $ ,平面 $ {PAD} \cap $ 平面 $ {ABCD} = {AD} $ ,
所以 $ {BG} \bot $ 平面 $ {PAD} $ .
        \item 连接 $ {PG} $ ,因为 $ \bigtriangleup {PAD} $ 为正三角形, $ G $ 为 $ {AD} $ 边的中点,
得 $ {PG}\bot {AD} $ ,由(1)知 $ {BG}\bot {AD} $ ,
$ {PG} \subset $ 平面 $ {PGB},{BG} \subset $ 平面 $ {PGB},{PG} \cap  {BG} = G $ ,
所以 $ {AD} \bot $ 平面 $ {PGB} $ ,因为 $ {PB} \subset $ 平面 $ {PGB} $ .
所以 $ {AD}\bot {PB} $ .
        \item 由(2)可得 $ {PB}\bot {AD},{BG}\bot {AD} $ ,
$ \because {AD}\parallel {BC} $ ,所以 $ {PB}\bot {BC},{BG}\bot {BC} $ ,
所以 $ \angle {PBG} $ 为二面角 $ A - {BC} - P $ 的平面角
因为 $ {PG} = {BG} = \dfrac{\sqrt{3}}{2}a $ ,所以 $ \angle {PBG} = {45}^{ \circ  } $ ；
        \item 当 $ F $ 为 $ {PC} $ 边的中点时,满足平面 $ {DEF}\bot $ 平面 $ {ABCD} $ ,证明如下：
取 $ {PC} $ 的中点 $ F $ ,连接 $ {DE} $ 、 $ {EF} $ 、 $ {DF} $ ,
在 $ \bigtriangleup  {PBC} $ 中, $ {FE}\left\|  {PB}\right\|  ,{FE} ⊄ $ 平面 $ {PGB},{PB} \subset $ 平面 $ {PGB} $
$ \therefore {FE} \parallel $ 平面 $ {PGB} $
在菱形 $ {ABCD} $ 中, $ {DG}\parallel {BE} $ 且 $ {DG} = {BE} $
$ {BEDG} $ 为平行四边形,则 $ {DE}\parallel {BG},{DE} ⊄ $ 平面 $ {PGB},{BG} \subset $ 平面 $ {PGB} $
$ \therefore {DE} \parallel $ 平面 $ {PGB} $
$ {EF} \cap  {DE} = E $ ,所以平面 $ {DEF}\square $ 平面 $ {PGB} $ ,
因为 $ {BG}\bot $ 平面 $ {PAD} $ ,所以 $ {BG}\bot {PG} $ ,
又因为 $ {PG}\bot {AD},{AD} \cap  {BG} = G $ ,
$ \therefore {PG} \bot $ 平面 $ {ABCD} $ ,而 $ {PG} \subset $ 平面 $ {PGB} $ ,
所以平面 $ {PGB} \bot $ 平面 $ {ABCD} $ ,
所以平面 $ {DEF}\bot $ 平面 $ {ABCD} $ .
    \end{enumerate}
\end{jieda}


\end{enumerate}


